Following \citet{hard2018gboard,xu2023federated}, we train one-layer LSTM LMs of $\sim6.4M$ parameters for mobile keyboard applications. These LMs are deployed on device to predict words during decoding time to facilitate user typing. 
We use next word prediction (NWP) accuracy on a hold-out set of devices to track the training progress, and also conduct A/B test in production, where following two metrics are reported:
\begin{enumerate*}[label=\color{purple}(\arabic*)]
    \item Word Modified Ratio (WMR), the ratio of words being modified during typing or after committed; improvement is shown by reduction;
    \item Word Per Minute (WPM): the number of committed words per minute.
\end{enumerate*}
LMs are trained for different language-locale combination in different populations. We study Spanish in Spain (es-ES), Indonesian in Indonesia (id-ID), Portuguese in Brazil (pt-BR) and Portuguese in Portugal (pt-PT). LMs are pre-trained on public multilingual C4 dataset~\citep{xue2020mt5} before private training with FL and DP. 

\paragraph{Algorithm Setting} We compare our BLT-DP-FTRL algorithm with \treeagg~\citep{kairouz21practical,xu2023federated} and \bandmf~\citep{choquette2023amplified} discussed in \cref{sec:background}. As far as we know, these two are the only DP algorithms actively used to train LMs in a production FL system. We follow the system and parameter configurations in~\citep{xu2023federated,choquette2023amplified,gboard_dp_blogpost} for baselines, and compare to \treeagg for pt-PT and id-ID, and \bandmf for es-ES and pt-BR. However, we highlight it is challenging to exactly reproduce the settings, especially for the min-sep parameter. The \bandmf algorithm are optimized for total round $n=2000$, band $\hat b = 400$ for es-ES, and $\hat b = 1000$ for pt-BR. We optimize \blt for total round $n=4000$\footnote{As the target round is usually less than 2000, $n=4000$ for \blt is less favorable compared to $n=2000$ used for \bandmf. \blt is robust to the target round $n$, and achieves stronger results with an inferior $n$.}, estimated min-sep $b=100$ for pt-PT, $b=400$ for es-ES and id-ID, and $b=1000$ for pt-BR. We use \ourblt for multi-participation and estimate the max-par based on $n/b$. For these $n,b$ settings, only $d=4$ buffers can achieve near-optimal loss in optimization, and \ourblt matrices are parameterized by only 8 numbers (these parameters are provided in \cref{app:blt_params}). Though different configures are used for populations with different sizes, the BLT parameters $(\theta, \outp) \in \R^8$ optimized for $b=400$ can achieve competitive results for a wide range of min-seps, which can be a reasonable default for BLT-DP-FTRL. As discussed in \cref{tab:mechanism_summary}, \blt is more memory-efficient than both \treeagg and \bandmf. In the simulation results \cref{tab:sonwp}, \blt is also better than \treeagg for privacy-utility trade-off, and comparable with \bandmf. The results in production further show that the flexibility of \blt makes it easier to use in practice, and achieve better results than both \treeagg and \bandmf.